\section{渐近距离处的数据分析策略}
\label{sec:regge}

对有限强子动量，准 LF 关联的格点计算最终只能给出有限 $\lambda_{\rm L}=P^zz_{\rm L}$ 处的数据；由于大 $z$ 处有限体积修正增大、信噪比变差，$z_{\rm L}$ 通常小于格点尺寸。然而，要重构完整的部分子分布，我们需要所有准 LF 距离处的关联。

在有限动量下，准 LF 关联一般具有有限的关联长度（在 $\MS$ 或混合方案下），并在大 $z$ 处呈指数衰减。这类似于密度-密度~\cite{Burkardt:1994pw} 或流-流关联的情形，因为在辅助场表述中准 LF 关联可视为两个重到轻流之积~\cite{Samuel:1978iy,Gervais:1979fv,Arefeva:1980zd,Dorn:1986dt}。因此，与仅因 Regge 行为而在大 $\lambda$ 代数衰减的 LF（twist-2）关联相比，其傅里叶变换在有限 $z$ 或 $\lambda$ 处收敛较快。若 $z_{\rm L}$ 足够大使准 LF 关联几乎降为零，则可以作截至 $z_{\rm L}$ 的截断傅里叶变换得到准 PDF，由此产生的系统不确定性相对其他来源可以忽略。

然而在实际的格点计算中，$z_L$ 的选取受限于准 LF 关联快速增长的误差；大强子动量时尤其如此。于是，当我们按目标误差选定某个 $z_L$ 或 $\lambda_{\rm L}$ 时，准 LF 关联在该点的取值可能仍然相当可观。此时截断傅里叶变换会在准 PDF 中引起非物理振荡和不准确的小 $x$ 结果，这可以被表述为一个反演问题~\cite{Karpie:2019eiq}。文献中已采用若干策略来处理这一问题，例如 Backus-Gilbert 方法~\cite{Karpie:2019eiq,Bhat:2020ktg}、神经网络与贝叶斯重构~\cite{Karpie:2019eiq}、压制长程关联的高斯重加权方法~\cite{Ishikawa:2019flg}、相当于作分部积分并忽略截断点边界项的导数方法~\cite{Lin:2017ani}，以及利用高斯过程回归在整个定义域上重构连续形式准 LF 关联的 Bayes-Gauss-Fourier 变换~\cite{Alexandrou:2020tqq}。然而这些策略所采用的假设大多基于数学而非物理的理由。本文提议利用准 LF 关联渐近行为的知识，进行具有物理动机的外推至 $\lambda\!=\!\infty$。
外推之后，可以对 $|z|\le z_{\rm L}$ 作离散傅里叶变换——其离散化误差可通过格距依赖加以研究——而对外推的部分则可解析地完成傅里叶变换。这样，数学上的反演问题便由物理考量解决。
虽然这种外推并不能提供小 $x$ PDF 的第一性预言，但它有助于消除非物理振荡，并为估计该区域的系统不确定性提供一种合理的途径。

根据动量的大小，我们提议外推采用指数或代数两种衰减形式之一。下面详细讨论这两种做法，并说明如何估计相应的不确定性。


\subsection{中等大 $P^z$ 下的指数外推}

如上所述，对中等大的动量 $P^z$，准 LF 关联一般具有有限的关联长度：$z$ 空间中为 $\xi_z \sim 1/\Lambda_{\rm QCD}$，
LF 距离 $\lambda$ 空间中为
$\xi_\lambda \sim P^z/\Lambda_{\rm QCD}$。这是非微扰 QCD 的禁闭性质以及该关联类空性质的后果。有限的关联长度伴随着指数衰减 $\exp(-z/\xi_z)$，后者在大 $z$ 处变得显著。在准 LF 关联表现出指数衰减行为之前，它由随 $P^z$ 缓慢演化的领头扭度贡献主导。在 $\lambda$ 空间中，不同 $P^z$ 下的准 LF 关联可由 \fig{extr} 定性地描述。随着 $P^z$ 增大，准 LF 关联愈发接近领头扭度贡献，并在更大的 $\lambda$ 值处才开始呈现指数衰减。在 $P^z\to\infty$ 极限下，$\xi_\lambda$ 趋于无穷，准 LF 关联只剩在大 $\lambda$ 代数衰减的领头扭度贡献。



\begin{figure}[t]
\centering
\includegraphics[width=\columnwidth]{images/extrapolation}
\caption{不同 $P^z$ 下准 LF 关联在 $\lambda$ 空间的定性行为。对有限 $P^z$，短 $\lambda$ 处的关联近似为领头扭度贡献并随 $P^z$ 缓慢演化；大 $\lambda$ 处关联开始呈现指数衰减行为，且起始点与关联长度 $\xi_\lambda$ 都随 $P^z$ 增大。在 $P^z\to\infty$ 极限下，$\xi_\lambda$ 趋于无穷，准 LF 关联只含代数衰减的领头扭度贡献。}
\label{fig:extr}
\end{figure}

因此，当 $P^z$ 不很大（例如质子情形约 2--5 GeV）时，我们提议用指数衰减形式 $\sim e^{-\lambda/\xi_\lambda}$ 作外推（也可以在其上再叠加一些代数行为，以更好地表示准 LF 关联的 $\lambda$ 依赖）。注意，要使外推可控，关键在于格点数据须在误差变得过大之前呈现出指数衰减。这应通过更大的格点体积和/或更高统计量的测量来实现，这对当代计算资源是可行的。理想情形下，在足够大的 $z$ 或 $\lambda$ 处，准 LF 关联会在目标误差内实际上变为零，外推除极小 $x$ 区域外几乎不影响最终结果。更现实的情形是：格点数据表现出指数式下降，但在 $\lambda_{\rm L}$ 处仍有统计上显著的非零值，此时我们可以把指数外推做到 $\lambda=\infty$。

注意，尽管指数形式有物理动机，应当把哪些 $z$ 或 $\lambda$ 的格点数据纳入拟合仍不清楚，因为大 $z$ 数据仍含有可能掩盖结果的领头扭度贡献。也就是说，不同的拟合区间选择可能拟合出不同的关联长度 $\xi_\lambda$。不过，$\xi_\lambda$ 的变化主要影响非常小的 $x$ 区域，那里本来就因幂次修正而预言能力较弱。因此并非必须精确拟合 $\xi_\lambda$。相反，应利用这一性质变动拟合区间（例如在 $z_{\rm L}-5a\le z \le z_{\rm L}$ 内），检验最终结果对不同 $\xi_{\lambda}$ 的稳定性。

最后同样重要的是：指数衰减关联的傅里叶变换总会在 $x=0$ 给出有限的准 PDF，这与 PDF 小 $x$ 处的 Regge 行为不同。此外，由于大 $x$（$x\to1$）处的 PDF 对代数衰减的长程 LF 关联同样敏感，准 PDF 在该区域也会偏离 PDF。这些都表明端点区域 $x\to0$ 与 $x\to1$ 幂次修正的重要性，也提示我们如何估计指数外推的系统不确定性。具体而言，可以对同一范围的数据作代数外推（见下一小节）——这本质上等价于忽略大 $z$ 处的全部幂次修正——并把其与指数外推之差取为误差。可以预期，随着 $x$ 趋向端点，这一估计将给出越来越大的系统误差，这与动量空间展开的精度相一致。

\subsection{超大 $P^z$ 下的代数外推}

当未来的格点资源使 $P^z$ 变得非常大时，
$\xi_\lambda$ 也变得非常大，需要在更大的 $\lambda$ 值才能在格点数据中看到指数衰减。我们预期当 $\lambda\sim\lambda_{\rm L}$ 时衰减行为更像代数规律而非指数规律。换言之，由于 $\lambda \le \lambda_{\rm L}$ 的幂次修正预期受到良好压制，准 LF 关联非常接近领头扭度关联。此时可以用代数形式外推到 $\lambda=\infty$。


领头扭度关联的代数衰减是其关联长度 $\xi_\lambda$ 无穷的后果，并与渐近 Regge 行为相关~\cite{Regge:1959mz}。众所周知，小 $x$ 处部分子分布的渐近行为如 Regge 理论所示，形如 $x^a$。对非单态组合，领头 Regge 轨迹表明 $a\sim-1/2$。对单态组合，演化下与胶子分布的混合使情况更为微妙。在所谓的软玻密子（soft pomeron）模型中 $a\sim -1$。然而，大动量转移的散射数据表明存在更奇异的渐近行为，反映可能需要硬玻密子（hard pomeron）的贡献~\cite{Devenish:2004pb}。大 $x$ 处的渐近行为由夸克计数规则决定~\cite{Brodsky:1973kr}。当 $x\to 1$ 时，强子动量被被打出的夸克带走，其他旁观部分子分不到动量。渐近行为预言为 $(1-x)^b$，其中 $b=2n_s-1+2|\Delta S^z|$，$n_s$ 是最少的旁观部分子数目，$\Delta S^z$ 是被打出部分子与母强子自旋投影之差~\cite{Devenish:2004pb,Brodsky:2005wx}。例如，对质子中的价夸克，若被打出夸克的螺旋度与质子平行（反平行），则 $n_s=2$、$|\Delta S^z|= 0\, (1)$，从而 $b=3\, (5)$；而对 π 介子，由于 $n_s=1$、$|\Delta S^z|=1/2$，有 $b=2$。上述特征已被广泛用于 PDF 全局拟合：参数化 PDF 使其在 $x\to 0$ 表现为 $x^a$、在 $x\to 1$ 表现为 $(1-x)^b$，并对大量多样的实验数据拟合幂次 $a, b$。
这类幂律行为在全局拟合中的作用已在 ~\cite{Ball:2016spl,Nocera:2014uea} 中得到详细考察。

傅里叶变换到坐标空间后，上述渐近行为意味着纵向空间的关联按 $\lambda^{-\alpha}$（$\alpha$ 为与 $a, b$ 相关的正数）代数衰减而非指数衰减，因而关联长度为无穷。近期对 LF 波函数的分析在变换到共轭坐标空间时也观察到了类似的代数衰减行为~\cite{Brodsky:1997de,Miller:2019ysh}。

为了看清渐近行为如何帮助大动量下准 LF 关联的外推，我们从如下包含 $x\sim 0, 1$ 行为的简单 PDF 形式出发：
\begin{equation}
     x^a (1-x)^b \ .
\end{equation}
其坐标空间矩阵元定义为
\begin{equation}
h(\lambda)=\int_0^1 dx \, e^{i x\lambda}x^a (1-x)^b,
\end{equation}
由此可知在大 $\lambda$ 处
\begin{align}\label{asympt_coord}
h(\lambda)\sim\frac{\Gamma(1+a)}{(-i|\lambda|)^{a+1}}+e^{i\lambda}\frac{\Gamma(1+b)}{(i|\lambda|)^{b+1}} \ ,
\end{align}
其实部（虚部）是 $\lambda$ 的偶（奇）函数，这保证部分子分布在动量空间中是实函数。因此共轭 LF 关联在大 $\lambda$ 处按 $\lambda^{-\alpha(a,b)}$ 行为变化，其中
\begin{align}
\alpha(a,b)=\text{\rm min}(a+1,b+1).
\end{align}
在我们感兴趣的大多数情形中，$\alpha(a,b)=a+1$。施加式~(\ref{matching}) 的匹配会把光锥关联转换为准 LF 关联，同时给渐近行为带来对数修正。在因子化成立的区域，这些修正在领头对数（LL）精度下可重求和为 $\sim \exp(\gamma \ln z^2\mu^2)=(\lambda^2\mu^2/(p^z)^2)^\gamma$，它把准 LF 关联的渐近行为修改为
\begin{align}\label{asymphtilde}
\tilde h(\lambda,z)\sim e^{\gamma \ln z^2\mu^2}\frac{1}{|\lambda|^{\alpha(a,b)}}\sim |\lambda|^{-\alpha(a,b)+2\gamma }  \,.
\end{align}
这为足够大 $P^z$ 下大 $\lambda$ 处的准 LF 关联提供了有用的近似，并与前文基于关联长度的讨论一致。
现在可用如下代数形式把准 LF 关联外推至无穷 $\lambda$（以 $\lambda>0$ 为例）
\begin{equation}\label{extrap}
\frac{c_1}{(-i \lambda)^{d_1}}+e^{i\lambda}\frac{c_2}{(i \lambda)^{d_2}},
\end{equation}
它容纳了式~(\ref{asympt_coord}) 中两种不同结构。参数 $c_i, d_i$ 可以像指数外推那样拟合。最后，可以利用这些参数的不确定度来估计外推的系统误差。





用上述外推策略补充格点数据后，我们期待最终的 PDF 结果免于非物理振荡，并更好地收敛到物理区域 $0<x<1$。
